\documentclass[11pt]{article}
\usepackage[a4paper,margin=0.8in]{geometry}
\usepackage{calc}
\usepackage{array,longtable,booktabs}
\usepackage{xcolor}
\usepackage{hyperref}
\usepackage{enumitem}
\usepackage{fancyvrb}
\usepackage{fvextra}
\usepackage{framed}
\usepackage{color}
\fvset{breaklines=true, breakanywhere=true}
\DefineVerbatimEnvironment{Highlighting}{Verbatim}{breaklines,breakanywhere,commandchars=\\\{\}}
\usepackage{color}
\usepackage{fancyvrb}
\newcommand{\VerbBar}{|}
\newcommand{\VERB}{\Verb[commandchars=\\\{\}]}
\DefineVerbatimEnvironment{Highlighting}{Verbatim}{commandchars=\\\{\}}
% Add ',fontsize=\small' for more characters per line
\usepackage{framed}
\definecolor{shadecolor}{RGB}{248,248,248}
\newenvironment{Shaded}{\begin{snugshade}}{\end{snugshade}}
\newcommand{\AlertTok}[1]{\textcolor[rgb]{0.94,0.16,0.16}{#1}}
\newcommand{\AnnotationTok}[1]{\textcolor[rgb]{0.56,0.35,0.01}{\textbf{\textit{#1}}}}
\newcommand{\AttributeTok}[1]{\textcolor[rgb]{0.13,0.29,0.53}{#1}}
\newcommand{\BaseNTok}[1]{\textcolor[rgb]{0.00,0.00,0.81}{#1}}
\newcommand{\BuiltInTok}[1]{#1}
\newcommand{\CharTok}[1]{\textcolor[rgb]{0.31,0.60,0.02}{#1}}
\newcommand{\CommentTok}[1]{\textcolor[rgb]{0.56,0.35,0.01}{\textit{#1}}}
\newcommand{\CommentVarTok}[1]{\textcolor[rgb]{0.56,0.35,0.01}{\textbf{\textit{#1}}}}
\newcommand{\ConstantTok}[1]{\textcolor[rgb]{0.56,0.35,0.01}{#1}}
\newcommand{\ControlFlowTok}[1]{\textcolor[rgb]{0.13,0.29,0.53}{\textbf{#1}}}
\newcommand{\DataTypeTok}[1]{\textcolor[rgb]{0.13,0.29,0.53}{#1}}
\newcommand{\DecValTok}[1]{\textcolor[rgb]{0.00,0.00,0.81}{#1}}
\newcommand{\DocumentationTok}[1]{\textcolor[rgb]{0.56,0.35,0.01}{\textbf{\textit{#1}}}}
\newcommand{\ErrorTok}[1]{\textcolor[rgb]{0.64,0.00,0.00}{\textbf{#1}}}
\newcommand{\ExtensionTok}[1]{#1}
\newcommand{\FloatTok}[1]{\textcolor[rgb]{0.00,0.00,0.81}{#1}}
\newcommand{\FunctionTok}[1]{\textcolor[rgb]{0.13,0.29,0.53}{\textbf{#1}}}
\newcommand{\ImportTok}[1]{#1}
\newcommand{\InformationTok}[1]{\textcolor[rgb]{0.56,0.35,0.01}{\textbf{\textit{#1}}}}
\newcommand{\KeywordTok}[1]{\textcolor[rgb]{0.13,0.29,0.53}{\textbf{#1}}}
\newcommand{\NormalTok}[1]{#1}
\newcommand{\OperatorTok}[1]{\textcolor[rgb]{0.81,0.36,0.00}{\textbf{#1}}}
\newcommand{\OtherTok}[1]{\textcolor[rgb]{0.56,0.35,0.01}{#1}}
\newcommand{\PreprocessorTok}[1]{\textcolor[rgb]{0.56,0.35,0.01}{\textit{#1}}}
\newcommand{\RegionMarkerTok}[1]{#1}
\newcommand{\SpecialCharTok}[1]{\textcolor[rgb]{0.81,0.36,0.00}{\textbf{#1}}}
\newcommand{\SpecialStringTok}[1]{\textcolor[rgb]{0.31,0.60,0.02}{#1}}
\newcommand{\StringTok}[1]{\textcolor[rgb]{0.31,0.60,0.02}{#1}}
\newcommand{\VariableTok}[1]{\textcolor[rgb]{0.00,0.00,0.00}{#1}}
\newcommand{\VerbatimStringTok}[1]{\textcolor[rgb]{0.31,0.60,0.02}{#1}}
\newcommand{\WarningTok}[1]{\textcolor[rgb]{0.56,0.35,0.01}{\textbf{\textit{#1}}}}
\setlist[itemize]{leftmargin=*}
\setlist[enumerate]{leftmargin=*}
\hypersetup{colorlinks=true,linkcolor=blue,urlcolor=blue}
\providecommand{\tightlist}{%
  \setlength{\itemsep}{0pt}\setlength{\parskip}{0pt}}
\title{R Workshop -- Day 5 -- Regression, ANOVA and Generalised Linear Models}
\author{Applied R for Statistical Teaching, Research and Reproducible Analysis}
\date{}
\begin{document}
\maketitle
\tableofcontents
\newpage
\hypertarget{r-workshop-day-5-regression-anova-and-generalised-linear-models}{%
\section{R Workshop -- Day 5 -- Regression, ANOVA and Generalised Linear
Models}\label{r-workshop-day-5-regression-anova-and-generalised-linear-models}}

\textbf{Session 5 of 7 \textbar{} Duration: 3 hours} \textbf{Programme:}
Applied R for Statistical Teaching, Research and Reproducible Analysis

\hypertarget{learning-objectives}{%
\subsection{Learning objectives}\label{learning-objectives}}

By the end of this session you will be able to:

\begin{enumerate}
\def\labelenumi{\arabic{enumi}.}
\tightlist
\item
  Translate a research question into R formula notation.
\item
  Fit and interpret a multiple linear regression model with
  \texttt{lm()}.
\item
  Include categorical predictors and correctly explain the reference
  category.
\item
  Fit and interpret one interaction term.
\item
  Check residuals, fitted values and influential observations.
\item
  Run and interpret a one-way ANOVA where appropriate.
\item
  Report coefficients, uncertainty and model limitations in plain
  language.
\end{enumerate}

\hypertarget{capstone-link}{%
\subsection{Capstone link}\label{capstone-link}}

Today you fit the \textbf{principal statistical model} for your
capstone: a regression of post-training score on pre-training score,
attendance and training track (with an interaction), directly answering
the Day 2 research question.

\hypertarget{segment-plan-3-hours}{%
\subsection{Segment plan (3 hours)}\label{segment-plan-3-hours}}

\begin{longtable}[]{@{}
  >{\raggedright\arraybackslash}p{(\columnwidth - 4\tabcolsep) * \real{0.3000}}
  >{\raggedleft\arraybackslash}p{(\columnwidth - 4\tabcolsep) * \real{0.4000}}
  >{\raggedright\arraybackslash}p{(\columnwidth - 4\tabcolsep) * \real{0.3000}}@{}}
\toprule\noalign{}
\begin{minipage}[b]{\linewidth}\raggedright
Segment
\end{minipage} & \begin{minipage}[b]{\linewidth}\raggedleft
Time
\end{minipage} & \begin{minipage}[b]{\linewidth}\raggedright
Purpose
\end{minipage} \\
\midrule\noalign{}
\endhead
\bottomrule\noalign{}
\endlastfoot
Opening and recap & 10 min & Recap Day 4 inference, preview today's
model \\
Concept bridge & 25 min & Formula notation, what a coefficient means \\
Guided coding & 60 min & Fit, extend and diagnose a regression model \\
Participant practice & 45 min & Fit, interpret and diagnose the capstone
model \\
Interpretation and reporting & 25 min & Report coefficients and
limitations \\
Evidence log and capstone update & 15 min & Save the diagnosed model
report \\
\end{longtable}

\begin{center}\rule{0.5\linewidth}{0.5pt}\end{center}

\hypertarget{part-1-concept-bridge-core-workshop-activity}{%
\subsection{Part 1 -- Concept bridge (Core workshop
activity)}\label{part-1-concept-bridge-core-workshop-activity}}

\hypertarget{formula-notation}{%
\subsubsection{Formula notation}\label{formula-notation}}

\texttt{lm(post\_score\ \textasciitilde{}\ pre\_score\ +\ attendance\ +\ training\_track,\ data\ =\ lecturers)}
reads as: model \texttt{post\_score} as a function of
\texttt{pre\_score}, \texttt{attendance} and \texttt{training\_track}.
The left-hand side is the outcome; the right-hand side lists predictors
separated by \texttt{+}. An interaction between two predictors is
written \texttt{a:b}, or \texttt{a*b} as shorthand for
\texttt{a\ +\ b\ +\ a:b}.

\hypertarget{what-a-coefficient-means}{%
\subsubsection{What a coefficient
means}\label{what-a-coefficient-means}}

For a numeric predictor, the coefficient is the expected change in the
outcome for a one-unit increase in that predictor, holding other
predictors in the model constant. For a categorical predictor, R picks
one level as the \textbf{reference category} (by default, the first
alphabetically) and reports every other level's coefficient as a
difference \emph{from that reference}, holding other predictors
constant.

\begin{center}\rule{0.5\linewidth}{0.5pt}\end{center}

\hypertarget{part-2-guided-coding-core-workshop-activity}{%
\subsection{Part 2 -- Guided coding (Core workshop
activity)}\label{part-2-guided-coding-core-workshop-activity}}

\begin{Shaded}
\begin{Highlighting}[]
\FunctionTok{library}\NormalTok{(tidyverse)}
\FunctionTok{library}\NormalTok{(broom)}

\NormalTok{lecturers }\OtherTok{\textless{}{-}} \FunctionTok{read\_csv}\NormalTok{(}\StringTok{"Data/processed/anchor\_dataset\_clean.csv"}\NormalTok{, }\AttributeTok{show\_col\_types =} \ConstantTok{FALSE}\NormalTok{) }\SpecialCharTok{\%\textgreater{}\%}
  \FunctionTok{mutate}\NormalTok{(}\AttributeTok{training\_track =} \FunctionTok{factor}\NormalTok{(training\_track, }\AttributeTok{levels =} \FunctionTok{c}\NormalTok{(}\StringTok{"Online"}\NormalTok{, }\StringTok{"In{-}Person"}\NormalTok{, }\StringTok{"Blended"}\NormalTok{)))}
\end{Highlighting}
\end{Shaded}

Setting the factor levels explicitly fixes the reference category at
``Online'' rather than leaving it to alphabetical order, and documents
the choice.

\hypertarget{multiple-linear-regression}{%
\subsubsection{2.1 Multiple linear
regression}\label{multiple-linear-regression}}

\begin{Shaded}
\begin{Highlighting}[]
\NormalTok{model\_main }\OtherTok{\textless{}{-}} \FunctionTok{lm}\NormalTok{(post\_score }\SpecialCharTok{\textasciitilde{}}\NormalTok{ pre\_score }\SpecialCharTok{+}\NormalTok{ attendance }\SpecialCharTok{+}\NormalTok{ training\_track, }\AttributeTok{data =}\NormalTok{ lecturers)}
\FunctionTok{summary}\NormalTok{(model\_main)}
\FunctionTok{tidy}\NormalTok{(model\_main, }\AttributeTok{conf.int =} \ConstantTok{TRUE}\NormalTok{)}
\end{Highlighting}
\end{Shaded}

\textbf{Reading the output:} the intercept is the predicted post-score
for an Online participant with \texttt{pre\_score\ =\ 0} and
\texttt{attendance\ =\ 0} -- rarely meaningful on its own, but necessary
algebraically. The \texttt{training\_trackIn-Person} and
\texttt{training\_trackBlended} coefficients are the expected score
difference for those tracks \emph{relative to Online}, holding pre-score
and attendance fixed.

\hypertarget{one-interaction}{%
\subsubsection{2.2 One interaction}\label{one-interaction}}

\begin{Shaded}
\begin{Highlighting}[]
\NormalTok{model\_interaction }\OtherTok{\textless{}{-}} \FunctionTok{lm}\NormalTok{(post\_score }\SpecialCharTok{\textasciitilde{}}\NormalTok{ pre\_score }\SpecialCharTok{+}\NormalTok{ attendance }\SpecialCharTok{*}\NormalTok{ training\_track, }\AttributeTok{data =}\NormalTok{ lecturers)}
\FunctionTok{summary}\NormalTok{(model\_interaction)}
\end{Highlighting}
\end{Shaded}

The interaction \texttt{attendance:training\_track} asks: does the
\emph{effect of attendance} on post-score differ by track? Compare
\texttt{model\_main} and \texttt{model\_interaction} with
\texttt{anova(model\_main,\ model\_interaction)} to see whether the
interaction earns its place in the model.

\begin{Shaded}
\begin{Highlighting}[]
\FunctionTok{anova}\NormalTok{(model\_main, model\_interaction)}
\end{Highlighting}
\end{Shaded}

\hypertarget{diagnostics}{%
\subsubsection{2.3 Diagnostics}\label{diagnostics}}

\begin{Shaded}
\begin{Highlighting}[]
\FunctionTok{par}\NormalTok{(}\AttributeTok{mfrow =} \FunctionTok{c}\NormalTok{(}\DecValTok{2}\NormalTok{, }\DecValTok{2}\NormalTok{))}
\FunctionTok{plot}\NormalTok{(model\_main)}
\FunctionTok{par}\NormalTok{(}\AttributeTok{mfrow =} \FunctionTok{c}\NormalTok{(}\DecValTok{1}\NormalTok{, }\DecValTok{1}\NormalTok{))}
\end{Highlighting}
\end{Shaded}

Read the four panels as: residuals vs fitted (look for curvature -- a
violation of linearity), Q-Q plot (look for departure from the line --
non-normal residuals), scale-location (look for a fan shape --
non-constant variance), and residuals vs leverage (look for points
outside Cook's distance contours -- influential observations).

\begin{Shaded}
\begin{Highlighting}[]
\NormalTok{influence\_check }\OtherTok{\textless{}{-}} \FunctionTok{augment}\NormalTok{(model\_main) }\SpecialCharTok{\%\textgreater{}\%}
  \FunctionTok{arrange}\NormalTok{(}\FunctionTok{desc}\NormalTok{(.cooksd)) }\SpecialCharTok{\%\textgreater{}\%}
  \FunctionTok{select}\NormalTok{(.rownames, .cooksd) }\SpecialCharTok{\%\textgreater{}\%}
  \FunctionTok{head}\NormalTok{(}\DecValTok{5}\NormalTok{)}
\NormalTok{influence\_check}
\end{Highlighting}
\end{Shaded}

\hypertarget{one-way-anova}{%
\subsubsection{2.4 One-way ANOVA}\label{one-way-anova}}

\begin{Shaded}
\begin{Highlighting}[]
\NormalTok{anova\_track }\OtherTok{\textless{}{-}} \FunctionTok{aov}\NormalTok{(post\_score }\SpecialCharTok{\textasciitilde{}}\NormalTok{ training\_track, }\AttributeTok{data =}\NormalTok{ lecturers)}
\FunctionTok{summary}\NormalTok{(anova\_track)}
\end{Highlighting}
\end{Shaded}

A one-way ANOVA on
\texttt{post\_score\ \textasciitilde{}\ training\_track} tests whether
mean post-score differs across the three tracks, ignoring
\texttt{pre\_score} and \texttt{attendance}. It answers a narrower
question than the regression model above, and is appropriate when you
genuinely only care about one categorical factor.

\hypertarget{reporting}{%
\subsubsection{2.5 Reporting}\label{reporting}}

\textbf{Reporting template:} ``Controlling for pre-training score and
attendance, {[}In-Person/Blended{]} participants scored on average
\texttt{coefficient} points {[}higher/lower{]} than Online participants
(95\% CI: \texttt{lower} to \texttt{upper}, p = \texttt{p\_value}). The
model explained \texttt{r.squared} of the variance in post-training
scores. {[}Note any diagnostic concern, e.g.~one influential observation
flagged by Cook's distance.{]}''

\begin{Shaded}
\begin{Highlighting}[]
\FunctionTok{glance}\NormalTok{(model\_main)}\SpecialCharTok{$}\NormalTok{r.squared}
\end{Highlighting}
\end{Shaded}

\begin{center}\rule{0.5\linewidth}{0.5pt}\end{center}

\hypertarget{part-3-participant-practice-core-workshop-activity}{%
\subsection{Part 3 -- Participant practice (Core workshop
activity)}\label{part-3-participant-practice-core-workshop-activity}}

\begin{enumerate}
\def\labelenumi{\arabic{enumi}.}
\tightlist
\item
  Fit \texttt{model\_main} and \texttt{model\_interaction} on your own
  copy of the cleaned anchor dataset.
\item
  Write, in plain language, what the \texttt{training\_track}
  coefficients mean relative to your chosen reference category.
\item
  Run the four-panel diagnostic plot and note, in a comment, anything
  that concerns you (or confirms the model is reasonable).
\item
  Identify the single most influential observation by Cook's distance
  and decide, with a one-sentence justification, whether it should be
  investigated further or left in the model.
\item
  Run the one-way ANOVA on \texttt{training\_track} alone and compare
  its conclusion to the regression's \texttt{training\_track}
  coefficients.
\item
  Save a short written model report
  (\texttt{Outputs/day5\_model\_report.md} or \texttt{.R} with comments)
  covering points 1-5.
\end{enumerate}

A facilitator-reviewed solution is in
\texttt{Solution-Scripts/day5\_regression\_solution.R}.

\begin{center}\rule{0.5\linewidth}{0.5pt}\end{center}

\hypertarget{optional-demonstration}{%
\subsection{Optional demonstration}\label{optional-demonstration}}

\begin{itemize}
\tightlist
\item
  \textbf{Logistic regression} -- model \texttt{completed} (Yes/No) with
  \texttt{glm(completed\_numeric\ \textasciitilde{}\ ...,\ family\ =\ binomial)}.
\item
  \textbf{Poisson regression} -- for count outcomes, if a suitable
  variable is available or simulated.
\item
  \textbf{Nested-model comparison} --
  \texttt{anova(model\_a,\ model\_b)} beyond the one interaction already
  shown.
\item
  \textbf{Advanced influence measures} -- DFBETAS, DFFITS via
  \texttt{car::influence.measures()}.
\item
  \textbf{Two-way ANOVA} --
  \texttt{aov(post\_score\ \textasciitilde{}\ training\_track\ *\ institution\_type,\ data\ =\ lecturers)},
  if time allows.
\end{itemize}

\begin{Shaded}
\begin{Highlighting}[]
\CommentTok{\# Optional: logistic regression on completion}
\NormalTok{lecturers\_glm }\OtherTok{\textless{}{-}}\NormalTok{ lecturers }\SpecialCharTok{\%\textgreater{}\%} \FunctionTok{mutate}\NormalTok{(}\AttributeTok{completed\_numeric =} \FunctionTok{if\_else}\NormalTok{(completed }\SpecialCharTok{==} \StringTok{"Yes"}\NormalTok{, }\DecValTok{1}\NormalTok{, }\DecValTok{0}\NormalTok{))}
\NormalTok{model\_logit }\OtherTok{\textless{}{-}} \FunctionTok{glm}\NormalTok{(completed\_numeric }\SpecialCharTok{\textasciitilde{}}\NormalTok{ attendance }\SpecialCharTok{+}\NormalTok{ satisfaction, }\AttributeTok{data =}\NormalTok{ lecturers\_glm, }\AttributeTok{family =}\NormalTok{ binomial)}
\FunctionTok{summary}\NormalTok{(model\_logit)}
\end{Highlighting}
\end{Shaded}

\begin{center}\rule{0.5\linewidth}{0.5pt}\end{center}

\hypertarget{reference-or-take-home-material}{%
\subsection{Reference or take-home
material}\label{reference-or-take-home-material}}

\begin{itemize}
\tightlist
\item
  Model-selection cautions: do not chase the highest R-squared by adding
  every available variable; justify each predictor from the research
  question.
\item
  Post-hoc comparison examples:
  \texttt{emmeans::emmeans(anova\_track,\ pairwise\ \textasciitilde{}\ training\_track)}
  for pairwise track comparisons after a significant ANOVA.
\end{itemize}

\begin{center}\rule{0.5\linewidth}{0.5pt}\end{center}

\hypertarget{daily-deliverable}{%
\subsection{Daily deliverable}\label{daily-deliverable}}

A diagnosed model report for the capstone dataset: the fitted model, its
coefficients with confidence intervals, a diagnostic-plot summary, and a
plain-language interpretation.

\hypertarget{evidence-log-entry}{%
\subsubsection{Evidence log entry}\label{evidence-log-entry}}

Complete the Day 5 row in \texttt{Workshop-Evidence-Log.md}: name your
principal model and note one limitation you would mention if presenting
this result to a non-statistician.

\begin{center}\rule{0.5\linewidth}{0.5pt}\end{center}

\hypertarget{facilitator-notes}{%
\subsection{Facilitator notes}\label{facilitator-notes}}

\begin{itemize}
\tightlist
\item
  By construction, the synthetic dataset has a real
  \texttt{training\_track} effect (Blended \textgreater{} In-Person
  \textgreater{} Online) and a real positive \texttt{attendance} effect,
  so a correctly specified \texttt{model\_main} should show both as
  statistically meaningful -- use this as your answer key.
\item
  Watch for participants treating the regression intercept as if it
  described a typical participant; it describes a hypothetical one with
  all numeric predictors at zero, which is rarely realistic here.
\item
  The Day 5 model becomes the Day 6 audit target -- make sure
  participants keep a copy of \texttt{model\_main}'s formula and
  coefficients accessible (e.g.~save the fitted object with
  \texttt{saveRDS()}), since Day 6 reproduces part of it by hand.
\end{itemize}
\end{document}
